\documentclass[class=llncs, crop=false]{standalone}
\input{preamble.tex}
\input{macros.tex}

\begin{document}
\section{Introduction}\label{sec:intro}
%
The area of automated reasoning known as
\emph{satisfiability modulo theories} (SMT)
focuses on the development of tools for
checking the satisfiability of first-order
formulas over combinations of mathematical
theories~\cite{Barrett2021}.
%
Such tools (known as SMT \emph{solvers}) are
increasingly used as back-ends for various tasks:
hardware~\cite{DBLP:conf/fmcad/MattareiMBDHH18} and
software verification~\cite{barnett06boogie,riant23debugging,kosmatov16frama-c},
%
model checking \cite{DBLP:conf/cav/ChampionMST16},
cyber-security \cite{Backes2018,rungta22billion},
and to improve automation in proof
assistants~\cite{DBLP:conf/cpp/ArmandFGKTW11,DBLP:journals/jar/BlanchetteBP13}.
%
Because these applications demand a high level
of confidence in the results produced by the solver,
many solvers are capable of generating
\emph{proof certificates}
that witness the unsatisfiability of formulas
reported as such by the solver.
%

% Towards the aim of standardizing and benchmarking solvers,
% the \emph{SMT library initiative} oversees the development of
The \emph{SMT-LIB standard}~\cite{Barrett2015-standard}
specifies the syntax and semantics of a
common language used for interacting with solvers.
%%
Proof generation is however not covered by the standard,
leading to the development of various proof formats,
namely LFSC~\cite{DBLP:journals/fmsd/StumpORHT13},
Alethe~\cite{schurrAletheGenericSMT2021},
and recently Eunoia~\cite{EthosUser_manualmdMain}.
%
Eunoia is both a proof output language and a logical framework
for specifying the constants, side-condition programs,
and inference rules of an SMT solver (e.g., \cvc{5}).
The \emph{Cooperating Proof Calculus} (CPC) is the
Eunoia signature encoding \cvc{5}'s proof system,
and Eunoia proof scripts may be checked by the C++ tool
Ethos~\cite{EthosUser_manualmdMain}.
%
%

The $\lambda\Pi$-calculus modulo rewriting~\cite{CousineauDowek}
and its implementations in the tools Dedukti~\cite{assaf2023dedukti}
and LambdaPi~\cite{hondet:hal-02981561} provide
trustworthy proof checking with an emphasis on
interoperability of proof systems.
%
This paper presents a translation from Eunoia to LambdaPi
and an OCaml tool $\texttt{eo2lp}$ that implements the
translation.
%
We used our tool to translate \emph{CPC-MINI}, a fragment
of CPC covering 14~SMT-LIB logic fragments
(\autoref{fig:pipeline}), and to further translate and
typecheck over 1{,}000~proof scripts produced by running
\cvc{5} on $\texttt{unsat}$ problems from those fragments.
% ---------------------------------------------------------
\begin{boxfigure}[t!]{fig:pipeline}
	{{Overview of the translation pipeline.
				\cvc{5} produces Eunoia proof scripts using rules
				from the CPC signature, which may be checked by Ethos.
				\texttt{eo2lp} translates both the signature and the proofs
				to LambdaPi, which independently verifies them.}}
	%
	\centering
	\alttext{Pipeline diagram with three vertical groups.
		Left: SMT-LIB fragments (QF\_UF, QF\_LIA, UFLIA, ALIA, QF\_NIA, UF, \ldots).
		Middle: Eunoia artefacts---a CPC signature (theories, rules, programs)
		and a Eunoia proof script (.prf.eo).
		Right: LambdaPi artefacts---a translated CPC signature plus Prelude,
		and a LambdaPi proof file (.prf.lp).
		Horizontal arrows: cvc5 maps SMT-LIB fragments to Eunoia proof scripts;
		eo2lp maps both the Eunoia signature and proof script to LambdaPi.
		Downward arrows: ethos checks the Eunoia bundle, lambdapi check verifies
		the LambdaPi bundle, each producing a tick-or-cross result.}{%
		\input{fig-pipeline.tex}}
\end{boxfigure}
% ---------------------------------------------------------

% \subsubsection{Contribution.}
% Practically,
% \texttt{eo2lp} (\autoref{sec:results}) allows
% using LambdaPi as an independent proof checker for
% \cvc{5}: it translates all 22 modules of CPC-MINI
% (359 rules, 96 programs, 57 constants) and verifies
% 97\% of 1{,}036 proof scripts across 14 SMT-LIB
% fragments.
% %
% Conceptually, the translation gives an interpretation
% of Eunoia in the $λΠ$-calculus in the absence of a
% model-theoretic semantics like those given in the
% SMT-LIB standard.
% %
% Eunoia's (dependent) types are interpreted as terms of a
% type universe~$\Set$, inference rules as axioms,
% and side-conditions as rewrite rules.

\subsubsection{Related Work.}
The reduction of SMT proof checking to type
checking originates with
LFSC~\cite{DBLP:journals/fmsd/StumpORHT13};
Eunoia inherits this approach and extends it with
higher-order features similar to those of the
speculative SMT-LIB~3
proposal~\cite{barrettSMTLIB3Bringing,smt3-proposal,barbosa17hosmt}.
%
Brown~\cite{Brown2021} compares the prospects of interpreting
the proposed type system within type-theoretic and set-theoretic
frameworks; our encoding realizes the former via the $λΠ$-calculus
modulo rewriting.

The Dedukti/LambdaPi ecosystem is an interoperability hub for
proof assistants (Rocq~\cite{boespflug12coqine},
Isabelle~\cite{isabelle_dedukti},
HOL~Light~\cite{assaf15translating},
Matita~\cite{krajono},
Lean~\cite{vaishnav25lean4less})
and automated provers
(iProverModulo~\cite{burel20first},
ArchSAT~\cite{bury:tel-02612985},
Vampire~\cite{komel2025case},
TPTP~\cite{DBLP:conf/flairs/SutcliffeBB25}).
%
Coltellacci et al.~\cite{coltellacciReconstructionSMTProofs2024a,coltellacci25frocos}
reconstruct cvc5 proofs in LambdaPi via hand-written soundness lemmas for
each rule of the Alethe~\cite{schurrAletheGenericSMT2021} proof format.
%
Similarly, \textsc{lean-smt}~\cite{mohamed25leansmt}
reconstructs Eunoia proofs in Lean 4 using lemmas
that encode a subset of CPC's rules.
%
Our work instead targets Eunoia at the framework
level to allow encoding arbitrary
Eunoia signatures and proof scripts
in a type-theoretic framework.


\subsubsection*{Outline}
%
\autoref{sec:eunoia} introduces the fragment of Eunoia
targeted by our translation and its elaboration step;
\autoref{sec:lambdapi} defines the translation into the
$λΠ$-calculus modulo rewriting uniformly across rules,
programs, and proof scripts;
\autoref{sec:results} describes \texttt{eo2lp} and
reports the benchmark results.

%
% SMT proof reconstruction is also used within
% proof assistants:
% Isabelle/HOL replays Alethe
% proofs~\cite{DBLP:conf/cade/SchurrFD21,lachnitt25isabelle},
% SMTCoq~\cite{ekici17smtcoq} integrates solvers
% into Rocq, and IsaRare~\cite{lachnitt24isarare}
% verifies CPC rewrite rules in Isabelle/HOL.


% ---------------------------------------------------------
% Placed here so it floats to the top of page 2
\begin{boxfigure}[t!]{fig:eo-syntax}
	{Eunoia syntax: terms, attributes, built-in symbols, commands, proof commands.}
	\alttext{Five stacked grammar tables.
		First: terms---symbols, literals, applications, or binder applications;
		literals as int(n), rat(n,m), or str(s); parameters and program cases.
		Second: parameter attributes (:implicit, :list) and constant attributes
		(right/left-associative with optional nil terminator, chainable, pairwise,
		binder, arg-list).
		Third: built-in symbols grouped by family---core (eo::eq, eo::requires,
		eo::quote, eo::var), boolean (eo::and, eo::or, eo::not, eo::xor),
		arithmetic (eo::add, eo::mul, eo::neg, eo::zdiv), and
		list (eo::nil, eo::cons, eo::list\_concat, eo::list\_len).
		Fourth: standard commands---declare-const, declare-parameterized-const,
		declare-consts, define (macro definition), program (with signature and
		cases), declare-rule (with optional premises/premise-list, args,
		requires, and conclusion attributes), include.
		Fifth: proof commands---assume, assume-push, step (with rule, optional
		premises, optional args), and step-pop.}{%
		\vspace{4pt}\centerline{$\eoTermSyntax$}
		\boxfigurehline
		\centerline{$\eoAttrSyntax$}
		\boxfigurehline
		\centerline{$\eoBuiltinFig$}
		\boxfigurehline
		\centerline{$\eoCommandSyntax$}
		\boxfigurehline
		\centerline{$\eoPrfCommandSyntax$}\vspace{4pt}}
\end{boxfigure}
% ---------------------------------------------------------

\end{document}
